% 大爆炸（综述）
% license CCBYSA3
% type Wiki

本文根据 CC-BY-SA 协议转载翻译自维基百科\href{https://en.wikipedia.org/wiki/Big_Bang}{相关文章}。

\begin{figure}[ht]
\centering
\includegraphics[width=8cm]{./figures/7a2c26b638cb807a.png}
\caption{} \label{fig_BBang_1}
\end{figure}
“大爆炸”（The Big Bang）是一种物理理论，用以描述宇宙如何从最初高密度、高温的状态中膨胀而来。基于大爆炸概念的各种宇宙学模型能够解释广泛的天文现象，[1][2][3] 包括轻元素的丰度、宇宙微波背景辐射（CMB），以及宇宙的大尺度结构。宇宙的均匀性——即所谓的视界与平坦性问题——通过“宇宙暴胀”得到解释：那是一段发生在宇宙最早期的加速膨胀阶段。对宇宙膨胀速率的精确测量将最初奇点的时间推算为约 137.87 ± 0.02 亿年前，这被认为是宇宙的年龄。大量实证证据强烈支持这一“大爆炸事件”，该理论如今已被广泛接受。[4]

根据已知物理定律，将这种宇宙膨胀过程反向外推，模型显示宇宙早期处于极端炽热且高密度的原始状态。物理学目前仍缺乏一种被广泛接受的理论来完整描述大爆炸最初阶段的条件。[5] 随着宇宙膨胀，它逐渐冷却到足以形成亚原子粒子，随后又形成原子。这些原初元素——主要是氢，辅以少量氦和锂——在暗物质的引力作用下聚集，最终形成早期的恒星与星系。对超新星红移的观测表明，宇宙的膨胀正在加速，这一现象被归因于一种被称为“暗能量”的概念。

“宇宙膨胀”的思想最早由物理学家亚历山大·弗里德曼（Alexander Friedmann）于 1922 年提出，他在数学上推导出了“弗里德曼方程”。[6][7][8][9] 宇宙膨胀的最早观测依据被称为“哈勃定律”（Hubble’s Law），由物理学家爱德文·哈勃（Edwin Hubble）于 1929 年发表，他发现星系正以与其距离成正比的速率远离地球。独立于弗里德曼的理论工作与哈勃的观测，物理学家乔治·勒梅特（Georges Lemaître）于 1931 年提出，宇宙起源于一个“原初原子”（primeval atom），从而引入了现代意义上的“大爆炸”概念。1964 年，人们发现了宇宙微波背景辐射。随后几年中，对该辐射的观测表明：它在天空各方向上高度均匀，且能量—强度曲线的形状都与“大爆炸模型”预测的早期高温高密度状态相一致。至 20 世纪 60 年代末，大多数宇宙学家已确信竞争性的“稳恒态宇宙”模型是错误的。[10]

然而，观测宇宙中仍存在若干现象尚未能由“大爆炸模型”充分解释。其中包括：物质与反物质数量不对称（即重子不对称性）、星系周围暗物质的具体性质，以及“暗能量”的起源。[11]

